\documentclass[10pt,a4paper]{article}
\usepackage{amsmath,amssymb,amsthm,graphicx}
\usepackage{amsmath,amssymb,amsthm,mathtools}
\usepackage{hyperref}
\usepackage{cleveref}
\newtheorem{theorem}{Theorem}
\newtheorem{fact}[theorem]{Fact}
\newtheorem{definition}[theorem]{Definition}
\newtheorem{lemma}[theorem]{Lemma}
\newtheorem{proposition}[theorem]{Proposition}
\newtheorem{corollary}[theorem]{Corollary}
\newtheorem{claim}[theorem]{Claim}
\newtheorem{remark}[theorem]{Remark}
\newtheorem{observation}[theorem]{Observation}
\DeclareMathOperator*\disc{disc}
\DeclareMathOperator*\err{err}
\DeclareMathOperator*\opt{opt}
\DeclareMathOperator*\Exp{{\bf E}}
\DeclareMathOperator*\Prob{{\bf Pr}}
\newcommand{\polylog}{\operatorname{polylog}}
\let\citep\cite
\newcommand{\bool}{\left\{0,1\right\}}
\newcommand{\boolpm}{\left\{-1,1\right\}}
\newcommand{\sgn}{\mathrm{sgn}}
\newcommand{\ip}{\mathrm{IP}}
\newcommand{\bigo}[1]{O\left(#1\right)}
\newcommand{\numberthis}{\addtocounter{equation}{1}\tag{\theequation}}
\newcommand{\prblm}[1]{{\uppercase{\rm#1}}}
\newcommand{\MCSP}{\prblm{MCSP}}
\newcommand{\MKTP}{\prblm{MKTP}}
\newcommand{\SG}{\prblm{SUM\text{-}GREATER}}
\newcommand{\MKtP}{\mathrm{MKtP}}
\newcommand{\field}[1]{\mathbb{F}_{#1}}
\newcommand{\CNF}{{\sf CNF}}
\newcommand{\DNF}{{\sf DNF}}
\newcommand{\AND}{{\sf AND}}
\newcommand{\OR}{{\sf OR}}
\newcommand{\PARITY}{{\sf PARITY}}
\newcommand{\FORMULA}{{\sf FORMULA}}
\newcommand{\ANY}{{\sf ANY}}
\newcommand{\XOR}{{\sf XOR}}
\newcommand{\LTF}{{\sf LTF}}
\newcommand{\PTF}{{\sf PTF}}
\newcommand{\SYM}{{\sf SYM}}
\newcommand{\GIP}{{\sf GIP}}
\newcommand{\TT}[1]{{\tt tt}\!\left(#1\right)}
\newcommand{\CC}[1]{{\tt CC}\!\left(#1\right)}
\newcommand{\EI}{\aleph}
\newcommand{\given}{\mid}
\newcommand{\uniform}{U}
\newcommand{\LH}[1]{#1^{\mathrm{L}}}
\newcommand{\RH}[1]{#1^{\mathrm{R}}}
\newcommand{\myapprox}[1]{\ensuremath{\stackrel{#1}{\approx}}}
\newcommand{\real}{\mathbb{R}}
\newcommand{\naturals}{\mathbb{N}}
\newcommand{\abs}[1]{\left|#1\right|}
\newcommand{\prn}[1]{\left(#1\right)}
\newcommand{\cprn}[1]{\!\left(#1\right)}
\newcommand{\sqbra}[1]{\left[#1\right]}
\newcommand{\csqbra}[1]{\!\left[#1\right]}
\newcommand{\brkts}[1]{\left\{#1\right\}}
\newcommand{\cbrkts}[1]{\!\left\{#1\right\}}
\newcommand{\boool}{\brkts{-1,1}}
\newcommand{\angles}[1]{\left\langle#1\right\rangle}
\newcommand{\bars}[1]{\left\|#1\right\|}
\newcommand{\cpoly}[1]{\textup{poly}\cprn{#1}}
\newcommand{\quasipoly}[1]{\textup{quasipoly}\prn{#1}}
\newcommand{\cquasipoly}[1]{\textup{quasipoly}\cprn{#1}}
\newcommand{\var}[1]{\textbf{Var}\csqbra{#1}}
\newcommand{\cov}[2]{\textbf{Cov}\csqbra{#1,#2}}
\newcommand{\weight}[2]{{\bf W}^{#2}\csqbra{#1}}
\newcommand{\bin}[1]{{\rm Bin}\prn{#1}}
\newcommand{\cbin}[1]{{\rm Bin}\cprn{#1}}
\newcommand{\pr}[1]{\Prob\sqbra{#1}}
\newcommand{\cpr}[1]{\Prob\csqbra{#1}}
\newcommand{\prt}[2]{\Prob_{#2}\sqbra{#1}}
\newcommand{\cprt}[2]{\Prob_{#2}\csqbra{#1}}
\newcommand{\expect}[1]{\Exp\sqbra{#1}}
\newcommand{\cexpect}[1]{\Exp\csqbra{#1}}
\newcommand{\expectt}[2]{\Exp_{#2}\sqbra{#1}}
\newcommand{\cexpectt}[2]{\Exp_{#2}\csqbra{#1}}
\newcommand{\ACMOD}{{\sf AC}^0\circ{\sf MOD}_2}
\newcommand{\DNFMOD}{{\sf DNF}\circ{\sf MOD}_2}
\newcommand{\ANDMOD}{{\sf AND}\circ{\sf MOD}_2}
\newcommand{\formulaxor}{{\sf FORMULA}\circ{\sf XOR}}
\newcommand{\formulasizexor}[1]{{\sf FORMULA}\csqbra{#1}\circ{\sf XOR}}
\newcommand{\DNFMODAND}{{\sf DNF}\circ{\sf MOD}_2\circ{\sf AND}_t}

% Unused native Set renewal omitted.

\DeclarePairedDelimiter\ceil{\lceil}{\rceil}
\DeclarePairedDelimiter\floor{\lfloor}{\rfloor}
% Duplicate paired-delimiter declaration omitted.
% Duplicate paired-delimiter declaration omitted.
\allowdisplaybreaks
\AtBeginDocument{\author{Valentine Kabanets, Sajin Koroth, Zhenjian Lu\\Dimitrios Myrisiotis, Igor C. Oliveira}}
\title{Average-case generalized-inner-product lower bounds for formulas with low-communication leaf gates}
\author{Valentine Kabanets, Sajin Koroth, Zhenjian Lu, Dimitrios Myrisiotis, Igor C. Oliveira}
\date{doi:10.4230/LIPIcs.CCC.2020.15; CC BY 3.0 Unported}
\begin{document}
\maketitle
\noindent\textit{Editorial note.} Editable transcription only of corresponding formalCCBY3.0 Unported content; nonexclusive nativev1 is private aid. EntireSection3 andAppendixA.1 proofs retained with requiredSection2 definitions/cited prior facts. Original author probabilities/inequalities/constants retained without mathematical repair or new assumptions. MainTheorem1Items2/3 are externaloriginal scope references, not imported or counted here. Author metadata unchanged; title-page author list split into two lines to fit. Native unsupported/unused publisher packages omitted, necessarymath macros explicit, no input/include/assets. All production/validationYupeng; root owns compilation.\par\medskip
\noindent\textit{Allowed standard prerequisites (editorial).}\begin{itemize}
\item DeMorgan formula approximation, multiparty randomized/deterministic communication and GIP communication lower bound. FormalSection2 model/approximation statements and citedBNS lowerbound; source-local correlation conversion and fullTal-formula lemma proofs inAppendixA.1 included
\end{itemize}\medskip
	\paragraph{$\mathsf{GIP}^k_n$.}
	
	The Generalized Inner Product function $\GIP^k_{n}\colon \{0,1\}^{n} \to \{0,1\}$ is defined as
	\[
	    \GIP_{n}^{k}\left(x^{(1)},x^{(2)},\dots,x^{(k)}\right) = \sum_{j=1}^{n/k} \bigwedge_{i=1}^{k} x_{j}^{(i)}~(\mathsf{mod}\;2),
	\]
	where $x^{(i)}\in \{0,1\}^{n/k}$ for each $i\in [k]$.

\setcounter{section}{1}
\setcounter{theorem}{7}
	\section{Preliminaries}\label{sec:preliminaries}
	
	\subsection{Notation}
	
	Let $n\in\mathbb{N}$; we denote $\brkts{1,\dots,n}$ by $[n]$, and denote by $\uniform_n$ the uniform distribution over $\bool^n$. We use $\widetilde{O}(\cdot)$ (and $\widetilde{\Omega}(\cdot)$)  to hide polylogarithmic factors. That is, for any $f\colon\mathbb{N}\to\mathbb{N}$, we have that $\widetilde{O}(f(n))=O\cprn{f(n)\cdot\polylog\cprn{f(n)}}$.
	
	In this paper, we will mainly use $\boolpm$ as the Boolean basis. In some parts of this paper, we will use the $\bool$ basis for the simplicity of the presentation. This will be specified in corresponding sections.
	
	\subsection{De Morgan formulas and extensions}
	
	\begin{definition}
	    An $n$-variate \emph{de Morgan formula} is a directed rooted tree; its non-leaf vertices (henceforth, \emph{internal gates}) take labels from $\brkts{{\rm AND},{\rm OR},{\rm NOT}}=\brkts{\land,\lor,\neg}$ and its leaves (henceforth, \emph{variable gates}) take labels from the set of variables $\brkts{x_1,\dots,x_n}$. Each internal gate has bounded in-degree (henceforth, \emph{fan-in}); the {\rm NOT} gate in particular has fan-in $1$ and every variable gate has fan-in $0$. The \emph{size} of a de Morgan formula is the number of its leaf gates.
%
	\end{definition}
	
	In this work, we denote by $\FORMULA[s]$ the class of Boolean functions computable by size-$s$ de Morgan formulas. Let $\mathcal{G}$ denote some class of Boolean functions; then, we denote by $\FORMULA[s]\circ\mathcal{G}$ the class of functions computable by some size-$s$ de Morgan formula where its leaves are labelled by functions in $\mathcal{G}$.
	
%
	
%
%
%
%
%
%
%
%
%
    
%
%
%
%
%
    
    \subsection{Approximating polynomials}
    
    \begin{definition}[Point-wise approximation]
        For a Boolean function $f\colon\boolpm^n\to\boolpm$, we say that the function $\tilde{f}\colon\boolpm^n\to\real$ $\varepsilon$-approximates $f$ if for every $z\in\boolpm^n$,
		\[
			\left| f(z)-\tilde{f}(z)\right| \leq \varepsilon.
		\]
    \end{definition}
    
    We will need the following powerful result for the approximating degree of de Morgan formulas.
    
    \begin{theorem}[\cite{Rei11b}, see also~\cite{ BNRdW07}]\label{thm:approx-polynomial}
		Let $s>0$ be an integer and $0<\varepsilon<1$. Any de Morgan formula $F\colon\boolpm^n\to\boolpm$ of size $s$ has a $\varepsilon$-approximating polynomial of degree $d=O(\sqrt{s}\cdot\log(1/\varepsilon))$. That is, there exists a degree-$d$ polynomial $p\colon\boolpm^n\to\real$ over the reals such that for every $z\in\boolpm^n$,
		\[
			\left| p(z)-F(z)\right| \leq \varepsilon.
		\]
	\end{theorem}
	
	Note that \Cref{thm:approx-polynomial} still holds if we use $\bool$ as the Boolean basis.
    
	\subsection{Communication complexity}
	
	We use standard definitions from communication complexity. In this paper we consider the standard two party model of Yao and its generalizations to multiparty setting. We denote deterministic communication complexity of a Boolean function by $D(f)$ in the two party setting. We refer to \citep{Kushilevitz-Nisan97} for 
	standard definitions from communication complexity.
	
%
%
	    
%
%
	
%
	
	\begin{definition}
	    Let $f\colon\bool^{n}\to\bool$ be a Boolean function. The \emph{communication matrix of $f$}, namely $M_f$, is a $2^{n/2}\times2^{n/2}$ matrix defined by $\prn{M_f}_{x,y}:=f(x,y)$.
	\end{definition}
	
	\begin{definition}
	    A \emph{rectangle} is a set of the form $A\times B$, for $A,B\subseteq\bool^n$. A \emph{monochromatic rectangle} is a rectangle $S$ such that for all pairs $(x,y)\in S$ the value $f(x,y)$ is the same.
    \end{definition}
	
	\begin{lemma}
	    Let $\Pi$ be a protocol that computes $f\colon\bool^{n}\to\bool$ with at most $D$ bits of communication. Then, $\Pi$ induces a partition of $M_f$ into at most $2^D$ monochromatic rectangles.
	\end{lemma}
	
	Given a protocol, its \emph{transcript} is the sequence of bits communicated.
	
	\begin{lemma}
	    For every transcript $z$ of some communication protocol, the set of inputs $(x,y)$ that generate $z$ is a rectangle.
	\end{lemma}
	
%
%
%
	
%
%
%

    Below, we recount the definitions of two multiparty communication models used in this work, namely the number-on-forehead and the number-in-hand models.
	
	\begin{definition}[``Number-on-forehead'' communication model; informal]
	    In the $k$-party ``number-on-forehead'' communication model, 
%
	    there are $k$ players and $k$ strings $x_1,\dots,x_k\in\bool^{n/k}$ and player $i$ gets all the strings except for $x_i$. The players are interested in computing a value $f\cprn{x_1,\dots,x_k}$, where $f\colon\bool^{n}\to\bool$ is some fixed function. 
%
	    We denote by $D^{(k)}(f)$ the number of bits that must be exchanged by the best possible number on forehead protocol solving $f$.
    \end{definition}
        
%
    
    We also use the following weaker communication model.
    
    \begin{definition}[``Number-in-hand'' communication model; informal]
	    In the $k$-party ``number-in-hand'' communication model,
%
	    there are $k$ players and $k$ strings $x_1,\dots,x_k\in\bool^{n/k}$ and player $i$ gets only $x_i$.
%
	    The players are interested in computing a value $f\cprn{x_1,\dots,x_k}$, where $f\colon\bool^{n}\to\bool$ is some fixed function.
%
	    We denote by $D^{(k\text{-}\mathsf{NIH})}(f)$ the number of bits that must be exchanged by the best possible communication protocol.
    \end{definition}
    
    Note that $D^{(k\text{-}\mathsf{NIH})}(f)\leq\prn{1-1/k}\cdot n+1$, for any $n$-variate Boolean function $f$, as if $k-1$ players write on the blackboard their string, then the player that did not reveal her input may compute $f\cprn{x_1,\dots,x_k}$ on her own and then publish it.
    
    For the communication models mentioned above, there are also bounded-error randomized versions, denoted by $R_\delta$, $R_\delta^{(k)}$, and $R^{(k\text{-}\mathsf{NIH})}_\delta$, respectively, where $0<\delta<1$ is an upper bound on the error probability of the protocol. In this setting, the players have access to some shared random string, say $r$, and the aforementioned error probability of the protocol is considered over the possible choices of $r$. Moreover, we require the error to be at most $\delta$ on each fixed choice of inputs.
    
    We can extend the definitions of the communication complexity measures, defined above, to classes of Boolean functions, in a natural way. That is, for any communication complexity measure $M\in\brkts{D,D^{(k)},D^{(k\text{-}\mathsf{NIH})},R_\delta,R_\delta^{(k)},R_\delta^{(k\text{-}\mathsf{NIH})}}$ and for any class of Boolean functions $\mathcal{G}$, we may define
    \[
        M\cprn{\mathcal{G}}:=\max_{g\in\mathcal{G}}M\cprn{g}.
    \]
	We note that throughout this paper, we denote by $n$ the number of input bits for the function regardless the communication models. In the $k$-party communication setting (either NOF or NIH), we assume without loss of generality that $n$ is divisible by $k$.
	

\setcounter{theorem}{18}
	\section{Lower bounds}\label{sec:lower-bounds}
	
	In this section, we prove an average-case lower bound for the generalized inner product function against $\FORMULA\circ\mathcal{G}$, where $\mathcal{G}$ is the set of functions that have low-cost randomized communication protocols in the number-on-forehead setting. This corresponds to Item 1 of Theorem 1 of the original publication. Items 2 and 3 rely on our PRG constructions, and the proofs are deferred to Section 4 of the original publication.
	

	\begin{theorem}\label{thm:lb-main}
		For any integer $k\geq 2$, $s>0$ and any class of functions $\mathcal{G}$, let $C\colon\boolpm^{n}\to\boolpm$ be a function in $\FORMULA[s]\circ\mathcal{G}$ such that
		\[
		    \Prob_{x\sim\boolpm^{n}}\left[C(x)=\GIP_{n}^{k}(x)\right]\geq 1/2+\varepsilon.
		\]
		Then
		\[
	  		s = \Omega\left( \frac{n^2}{k^2\cdot16^{k}\cdot\left(R^{(k)}_{\varepsilon/(2n^2)}(\mathcal{G})+\log n\right)^2\cdot\log^2(1/\varepsilon)}\right).
		\]
	\end{theorem}

	
	We need a couple useful lemmas from~\cite{Tal16}, whose proofs are presented in \Cref{subsec:formula-lemmas} (\Cref{lem:lb-Tal-1-app} and \Cref{lem:lb-Tal-2-app}) for completeness.
	
	\begin{lemma}[\cite{Tal16}]\label{lem:lb-Tal-1}
	    Let $\mathcal{D}$ be a distribution over $\boolpm^n$, and let $f,C\colon\boolpm^n\to\boolpm$ be such that
	    \[
	        \Prob_{x\sim\mathcal{D}}[C(x)=f(x)]\geq 1/2+\varepsilon.
	    \] 
	    Let $\tilde{C}\colon\boolpm^n\to\real$ be a $\varepsilon$-approximating function of $C$, i.e., for every $x\in\boolpm^n$, $|C(x)-\tilde{C}(x)|\leq \varepsilon$. Then,
	    \[
	        \Exp_{x\sim\mathcal{D}}[\tilde{C}(x)\cdot f(x)]\geq \varepsilon.
	    \]
	\end{lemma}
	
	\begin{lemma}[\cite{Tal16}]\label{lem:lb-Tal-2}
	    Let $\mathcal{D}$ be a distribution over $\boolpm^n$ and let $\mathcal{G}$ be a class of functions. For $f\colon\boolpm^n\to\boolpm$, suppose that $D\colon\boolpm^n\to\boolpm\in \FORMULA[s]\circ\mathcal{G}$ is such that 
	    \[
	        \Prob_{x\sim\mathcal{D}}[D(x)=f(x)]\geq 1/2+\varepsilon_0.
	    \]
	    Then there exists some $h\colon\boolpm^n\to\boolpm\in \XOR_{O\left(\sqrt{s}\cdot\log(1/\varepsilon_0)\right)}\circ \mathcal{G}$ such that
	    \[
	        \Exp_{x\sim\mathcal{D}}[h(x)\cdot f(x)]\geq \frac{1}{s^{O\left(\sqrt{s}\cdot\log(1/\varepsilon_0)\right)}}.
	    \]
	\end{lemma}
    
    We also need the following communication-complexity lower bound for $\GIP$.
    
    \begin{theorem}[{\cite[Theorem 2]{BNS92}}]\label{thm:lb-GIP-bound}
        For any $k\geq 2$, any function that computes $\GIP_{n}^{k}$ on more than $1/2+\delta$ fraction of the inputs (over uniformly random inputs) must have $k$-party deterministic communication complexity at least $\Omega \! \left(n/(k \cdot 4^k) - \log(1/\delta)\right)$.
    \end{theorem}
    
    We first show that if a function correlates with some small formula, whose leaves are functions with low \emph{randomized} communication complexity, then it also correlates non-trivially with some function of relatively low \emph{deterministic} communication complexity.
    
    \begin{lemma}\label{thm:lb-proof-1}
        For any distribution $\mathcal{D}$ over $\boolpm^n$, and any class of functions $\mathcal{G}$, let $f\colon\boolpm^n\to\boolpm$ and $C\colon\boolpm^n\to\boolpm\in \FORMULA[s]\circ\mathcal{G}$ be such that 
	    \[
	        \Prob_{x\sim\mathcal{D}}[C(x)=f(x)]\geq 1/2+\varepsilon.
	    \]
	    Then there exists a function $h$, with $k$-party deterministic communication complexity at most
	    \[
	        O\left(R^{(k)}_{\varepsilon/(2s)}(\mathcal{G})\cdot\sqrt{s}\cdot\log(1/\varepsilon)\right),
	    \]
	    such that
	    \[
	        \Prob_{x\sim\mathcal{D}}[h(x)=f(x)]\geq 1/2+1/s^{O(\sqrt{s}\cdot\log(1/\varepsilon))}.
	    \]
    \end{lemma}
    
    \begin{proof}
        Let $C=F(g_1,g_2\dots,g_s)$ be the function in $\FORMULA[s]\circ\mathcal{G}$, where $F$ is a formula and $g_1,g_2,\dots,g_s$ are leaf functions from the class $\mathcal{G}$. For each  $g_i$, consider a $k$-party randomized protocol $ \Pi_i$ of cost at most $R=R^{(k)}_{\varepsilon/(2s)}(\mathcal{G})$ that has an error $\varepsilon/(2s)$.
        Now consider the following function
        \[
            \tilde{C}(x) \vcentcolon=\Exp_{\Pi_1,\Pi_2,\dots,\Pi_s}\left[D(x)\right],
        \]
        where
        \[
            D(x) \vcentcolon= F(\Pi_1(x),\Pi_2(x),\dots,\Pi_s(x)).
        \]
        Note that for any fixed choice of $(\Pi_1,\Pi_2,\dots,\Pi_s)$, $D$ is a formula whose leaves are functions with \emph{deterministic} communication complexity at most $R$. Next, we show the following.
        
        \begin{claim}\label{claim:lb-approximating}
            The function $\tilde{C}$ $\varepsilon$-approximates $C$.
        \end{claim}
        
        \begin{proof}[Proof of \Cref{claim:lb-approximating}]
            First note that since each $\Pi_i$ is a $(\varepsilon/(2s))$-error randomized protocol, by taking the union bound over the $s$ leaf functions, we have that for every input $x\in\boolpm^n$,
            \[
                \Prob_{\Pi_1,\Pi_2,\dots,\Pi_s}[\Pi_1(x)=g_1(x) \land \Pi_2(x)=g_2(x) \land\dots\land\Pi_s(x)=g_s(x)]\geq 1-\varepsilon/2.
            \]
            Denote by $\mathcal{E}$ the event $\Pi_1(x)=g_1(x) \land \Pi_2(x)=g_2(x)\land\dots\land\Pi_s(x)=g_s(x)$. We have for every $x\in\boolpm^n$,
            \begin{align*}
                \tilde{C}(x) &= \Exp_{\Pi_1,\Pi_2,\dots,\Pi_s}\left[D(x)\right]\\
                &=\Exp\left[D(x)\given \mathcal{E}\right]\cdot\Prob[\mathcal{E}]+\Exp\left[D(x)\given \neg\mathcal{E}\right]\cdot\Prob[\neg\mathcal{E}]\\
                &=C(x)\cdot\Prob[\mathcal{E}]+\Exp\left[D(x)\given \neg\mathcal{E}\right]\cdot\Prob[\neg\mathcal{E}].
            \end{align*}
            On the one hand, we have
            \[
                \tilde{C}(x)=C(x)\cdot\Prob[\mathcal{E}]+\Exp\left[D(x)\given \neg\mathcal{E}\right]\cdot\Prob[\neg\mathcal{E}]\leq  C(x) + \varepsilon/2.
            \]
            On the other hand, we get
            \[
                \tilde{C}(x)=C(x)\cdot\Prob[\mathcal{E}]+\Exp\left[D(x)\given \neg\mathcal{E}\right]\cdot\Prob[\neg\mathcal{E}]\geq C(x)\cdot(1-\varepsilon/2)+(-1)\cdot(\varepsilon/2)\geq C(x)-\varepsilon.
            \]
            This completes the proof of the claim.
        \end{proof}
        
        Now by \Cref{claim:lb-approximating} and \Cref{lem:lb-Tal-1}, we have
        \begin{equation}\label{eq:lb-eq-1}
	        \Exp_{x\sim\mathcal{D}}[\tilde{C}(x)\cdot f(x)]\geq \varepsilon.
        \end{equation}
        By the definition of $\tilde{C}$, \Cref{eq:lb-eq-1} implies that there exists some $D$, which is a formula whose leaves are functions with \emph{deterministic} communication complexity at most $R$, such that
        \[
            \Exp_{x\sim\mathcal{D}}[D(x)\cdot f(x)]\geq \varepsilon,
        \]
        which implies
        \[
            \Prob_{x\sim\mathcal{D}}[D(x)=f(x)]\geq 1/2+ \varepsilon/2.
        \]
        Then by \Cref{lem:lb-Tal-2}, there exists a function $h$, which can be expressed as the $\XOR$ of at most $O(\sqrt{s}\cdot\log(1/\varepsilon))$ leaf functions in $D$, such that
        \[
            \Exp_{x\sim\mathcal{D}}[h(x)\cdot f(x)]\geq \frac{1}{s^{O(\sqrt{s}\cdot\log(1/\varepsilon))}},
        \]
        which again implies
        \[
            \Prob_{x\sim\mathcal{D}}[h(x)=f(x)]\geq \frac{1}{2} + \frac{1}{s^{O\left(\sqrt{s}\cdot\log(1/\varepsilon)\right)}}.
        \]
        Finally, note that the $k$-party deterministic communication complexity of $h$ is at most 
        \[
            O(R\cdot\sqrt{s}\cdot\log(1/\varepsilon)),
        \]
        where $R=R^{(k)}_{\varepsilon/(2s)}(\mathcal{G})$.
    \end{proof}
    
    We are now ready to show \Cref{thm:lb-main}.
    
    \begin{proof}[Proof of \Cref{thm:lb-main}]
        Consider \Cref{thm:lb-proof-1} with $f$ being $\GIP_{n}^{k}$ and $\mathcal{D}$ being the uniform distribution. Consider \Cref{thm:lb-GIP-bound} with $\delta=1/s^{O(\sqrt{s}\cdot\log(1/\varepsilon))}$. We have
        \[
            O\left(R^{(k)}_{\varepsilon/(2s)}(\mathcal{G})\cdot\sqrt{s}\cdot\log(1/\varepsilon)\right)\geq n/(k4^k) - O\left(\sqrt{s}\cdot\log(s)\cdot\log(1/\varepsilon))\right),
        \]
        which implies
        \[
            s\geq \Omega\left( \frac{n^2}{k^2\cdot 16^{k}\cdot\left(R^{(k)}_{\varepsilon/(2n^2)}(\mathcal{G})+\log n\right)^2\cdot\log^2(1/\varepsilon)}\right).\qedhere
        \]
    \end{proof}
	

\appendix
\setcounter{theorem}{49}
    \section{Proofs of useful lemmas}
    
    \subsection{Useful lemmas for formulas}\label{subsec:formula-lemmas}
    
    The proofs in this section are essentially the same as that of~\cite{Tal16}.
    
    \begin{lemma}[\cite{Tal16}, \Cref{lem:lb-Tal-1} restated]\label{lem:lb-Tal-1-app}
	    Let $\mathcal{D}$ be a distribution over $\boolpm^n$, and let $f,C\colon\boolpm^n\to\boolpm$ be such that
	    \[
	        \Prob_{x\sim\mathcal{D}}[C(x)=f(x)]\geq 1/2+\varepsilon.
	    \] 
	    Let $\tilde{C}\colon\boolpm^n\to\real$ be a $\varepsilon$-approximating function of $C$, i.e., for every $x\in\boolpm^n$, $|C(x)-\tilde{C}(x)|\leq \varepsilon$. Then,
	    \[
	        \Exp_{x\sim\mathcal{D}}[\tilde{C}(x)\cdot f(x)]\geq \varepsilon.
	    \]
	\end{lemma}
	
	\begin{proof}
		Note that since $\tilde{C}$ $\varepsilon$-approximate $C$, we have for every $x\in\boolpm^n$
		\[
			\tilde{C}\cdot C(x)\geq 1-\varepsilon,
		\]
		and 
		\[
			\tilde{C}\cdot (1-C(x))\geq -1-\varepsilon.
		\]
		Then,
		\begin{align*}
			\Exp_{x\sim\mathcal{D}}[\tilde{C}(x)\cdot f(x)]&=\Exp_{x\sim\mathcal{D}}[\tilde{C}(x)\cdot f(x)\given C(x)=f(x)]\cdot\Prob_{x\sim\mathcal{D}}[C(x)=f(x)]\\
			&\qquad+\Exp_{x\sim\mathcal{D}}[\tilde{C}(x)\cdot f(x)\given C(x)\ne f(x)]\cdot\Prob_{x\sim\mathcal{D}}[C(x)\ne f(x)]\\
			&\geq(1-\varepsilon)\cdot \Prob_{x\sim\mathcal{D}}[C(x)=f(x)] + (-1-\varepsilon)\cdot \left(1-\Prob_{x\sim\mathcal{D}}[C(x)=f(x)]\right)\\
			&=2\cdot \Prob_{x\sim\mathcal{D}}[C(x)=f(x)] - 1 -\varepsilon\\
			&\geq 2\cdot (1/2+\varepsilon) - 1 - \varepsilon \geq \varepsilon,
		\end{align*}
		as desired.
	\end{proof}
	
	\begin{lemma}[\cite{Tal16}, \Cref{lem:lb-Tal-2} restated]\label{lem:lb-Tal-2-app}
	    Let $\mathcal{D}$ be a distribution over $\boolpm^n$ and let $\mathcal{G}$ be a class of functions. For $f\colon\boolpm^n\to\boolpm$, suppose that $D\colon\boolpm^n\to\boolpm\in \FORMULA[s]\circ\mathcal{G}$ is such that 
	    \[
	        \Prob_{x\sim\mathcal{D}}[D(x)=f(x)]\geq 1/2+\varepsilon_0.
	    \]
	    Then there exists some $h\colon\boolpm^n\to\boolpm\in \XOR_{O\left(\sqrt{s}\cdot\log(1/\varepsilon_0)\right)}\circ \mathcal{G}$ such that
	    \[
	        \Exp_{x\sim\mathcal{D}}[h(x)\cdot f(x)]\geq \frac{1}{s^{O\left(\sqrt{s}\cdot\log(1/\varepsilon_0)\right)}}.
	    \]
	\end{lemma}
	
    \begin{proof}
		Let 
		\[
			D = F(g_1,g_2\dots,g_s)
		\]
		be a device in $\FORMULA\circ\mathcal{G}$ where $F$ is a formula and $g_1,g_2,\dots,g_s$ are function from $\mathcal{G}$.

		Let $p\colon\boolpm^s\to\real$ be a $\varepsilon_0$-approximating polynomial for $F$ of degree $d=O(\sqrt{s}\cdot\log(1/\varepsilon_0))$. Note that we can write
		\[
			p(z)=\sum_{\substack{S\subseteq[s]: \\ |S|\leq d}}\hat{p}(S)\cdot\prod_{i\in S}z_i.
		\]
		Also, for each $S\subseteq[s]$, we have
		\[
			|\hat{p}(S)|=\left| \Exp_{z\in\boolpm^s}[p(z)\cdot\prod_{i\in S}z_i] \right| \leq  1+\varepsilon_0.
		\]
		Now let
		\[
			\tilde{D} \vcentcolon= p(g_1,g_2\dots,g_s).
		\]
		Note that $\tilde{D}$ is a $\varepsilon_0$-approximating function for $D$. Therefore, by \Cref{lem:lb-Tal-1-app}, we have
		\begin{align*}
			\varepsilon_0 &\leq \Exp_{x\sim\mathcal{D}}[D(x)\cdot f(x)]\\
			&= \Exp_{x\sim\mathcal{D}}\left[\left(\sum_{\substack{S\subseteq[s]: \\ |S|\leq d}}\hat{p}(S)\cdot\prod_{i\in S}g_i \right) \cdot f(x)\right]\\
			&=\sum_{\substack{S\subseteq[s]: \\ |S|\leq d}}\hat{p}(S)\cdot \Exp_{x\sim\mathcal{D}}\left[\prod_{i\in S}g_i\cdot f(x)\right]\\
			&\leq \sum_{\substack{S\subseteq[s]: \\ |S|\leq d}} (1+\varepsilon_0) \cdot \left| \Exp_{x\sim\mathcal{D}}\left[\prod_{i\in S}g_i\cdot f(x)\right] \right|.
		\end{align*}
		The above equation is the sum of at most $s^{O(d)}$ summands. Therefore, there exists some $S\subseteq[s]$ such that
		\[
			\left|\Exp_{x\sim\mathcal{D}}\left[\prod_{i\in S}g_i\cdot f(x)\right]\right| \geq \frac{\varepsilon_0}{(1+\varepsilon_0)\cdot s^{O(d)}}\geq \frac{1}{s^{O\left(\sqrt{s}\cdot\log(1/\varepsilon_0)\right)}},
		\]
		which implies that there exists some $h$, such that either $h=\prod_{i\in S}g_i$ or $h=-\prod_{i\in S}g_i$, and
		\[
			\Exp_{x\sim\mathcal{D}}\left[h(x)\cdot f(x)\right] \geq \frac{1}{s^{O\left(\sqrt{s}\cdot\log(1/\varepsilon_0)\right)}}.
		\]
		Finally, note that such $h$ can be expressed as the $\XOR$ of at most $d$ functions from $\mathcal{G}$.
    \end{proof}
    

\begin{thebibliography}{99}
\bibitem{BNRdW07}
Harry Buhrman, Ilan Newman, Hein R{\"{o}}hrig, and Ronald de~Wolf.
\newblock Robust polynomials and quantum algorithms.
\newblock {\em Theory Comput. Syst.}, 40(4):379--395, 2007.
\newblock
  \href{https://doi.org/10.1007/s00224-006-1313-z}{doi.org/10.1007/s00224-006-1313-z}.

\bibitem{BNS92}
L{\'{a}}szl{\'{o}} Babai, Noam Nisan, and Mario Szegedy.
\newblock Multiparty protocols, pseudorandom generators for logspace, and
  time-space trade-offs.
\newblock {\em J. Comput. Syst. Sci.}, 45(2):204--232, 1992.
\newblock
  \href{https://doi.org/10.1016/0022-0000(92)90047-M}{doi.org/10.1016/0022-0000(92)90047-M}.

\bibitem{Kushilevitz-Nisan97}
Eyal Kushilevitz and Noam Nisan.
\newblock {\em Communication Complexity}.
\newblock Cambridge University Press, 1997.

\bibitem{Rei11b}
Ben Reichardt.
\newblock Reflections for quantum query algorithms.
\newblock In {\em {SODA}}, pages 560--569, 2011.
\newblock \href{https://arxiv.org/abs/1005.1601}{arXiv:1005.1601}.

\bibitem{Tal16}
Avishay Tal.
\newblock The bipartite formula complexity of inner-product is quadratic.
\newblock {\em Electronic Colloquium on Computational Complexity {(ECCC)}},
  23:181, 2016.
\newblock \href{http://eccc.hpi-web.de/report/2016/181}{ECCC:TR16-181}.

\end{thebibliography}
\end{document}
